\documentclass[11pt]{article}

\usepackage{amssymb,bm,mathtools} % fonts
\usepackage{color,fullpage,graphicx,hyperref,parskip} % layout + media

\begin{document}

\title{Optimal charging of a battery}
\author{Manu Mannattil}
\maketitle

\section{Problem formulation}

The battery state at a given time $t$ is quantified by its state of charge (SOC) $x(t)$.
If the current passing through the battery during galvanostatic cycling is $u(t)$, then the SOC dynamics is given by
%
\begin{equation}
  \frac{d\, x(t)}{d\, t} = \frac{u(t)}{Q}.
\end{equation}
%
Here $Q > 0$ is the maximum amount of charge that the battery can hold.
Charging (discharging) occurs till time $T$ such that
%
\begin{equation}
  \int_{0}^{T} u(t) = \pm Q.
\end{equation}
%
Hence, during normal working conditions of the battery, the SOC satisfies the constraint $0 \leq x(t) \leq 1$.

The cost function to be minimized is the total power loss
%
\begin{equation}
  \mathcal{E}[x(t), u(t)] = \int_{0}^{T}dt\, R(x, t) u^{2},
  \quad
  \text{s.t.}
  \quad
  \left\{
  \begin{aligned}
    \dot{x} &= Q^{-1}u,\\
    x(0) &= 0,\\
    x(T) &= 1.
  \end{aligned}
  \right.
\end{equation}
%
Here $R(x, t)$ is the internal resistance of the battery, which is taken to be a function of the SOC and time.
We take it to be a general, nonlinear function that is independent of the current $u$.%
\footnote{Any dependence on the current can be converted to a dependence on $\dot{x}$.}

We form the Hamiltonian using Pontryagin's minimum principle
%
\begin{equation}
  H = R u^{2} + \lambda \frac{u}{Q}.
\end{equation}
%
In terms of the costate $\lambda$, the optimal control $u^{*}$ is easy to find
%
\begin{equation}
  \frac{dH}{du} = 0
  \quad
  \implies
  \quad
  u^{*} = -\frac{\lambda}{2RQ}.
\end{equation}
%
Hamilton's equations then read
%
\begin{equation}
%
\begin{aligned}
  \dot{x} &= -\frac{\lambda}{2RQ^{2}},\\
  \dot{\lambda} &= -\frac{\partial_{x}R\lambda^{2}}{4R^{2}Q^{2}}
\end{aligned}
%
\quad
\text{s.t.}
\quad
\left\{
\begin{aligned}
  x(0) &= 0,\\
  x(T) &= 1.
\end{aligned}
\right.
\end{equation}
%
To solve this TPBVP, we have to find $\lambda(0)$ such that $x(T) = 1$ (e.g., using the shooting method).

\end{document}
